% ============================================================
% donghx / PyQCD 4150 复现与对照报告
% 编译：XeLaTeX 两遍；16:9；不把缺失的外部产物写成通过
% ============================================================
\documentclass[UTF8,aspectratio=169,11pt]{ctexbeamer}

\usepackage{amsmath,amssymb}
\usepackage{booktabs}
\usepackage{tabularx}
\usepackage{array}
\usepackage{tikz}
\usetikzlibrary{arrows.meta,positioning,calc}
\usepackage{listings}
\usepackage{xcolor}
\usepackage{hyperref}

\definecolor{primary}{HTML}{17365D}
\definecolor{accent}{HTML}{1F77B4}
\definecolor{evidence}{HTML}{1E8449}
\definecolor{warning}{HTML}{C55A11}
\definecolor{muted}{HTML}{5B6573}
\definecolor{panel}{HTML}{F4F7FA}
\definecolor{good}{HTML}{EAF5EE}
\definecolor{warnbg}{HTML}{FFF4E8}

\setbeamersize{text margin left=0.55cm,text margin right=0.55cm}
\setlength{\emergencystretch}{2em}
\setbeamertemplate{navigation symbols}{}
\setbeamercolor{normal text}{fg=black!86,bg=white}
\setbeamercolor{structure}{fg=primary}
\setbeamercolor{frametitle}{fg=primary,bg=white}
\setbeamercolor{block title}{fg=primary,bg=panel}
\setbeamercolor{block body}{fg=black!86,bg=panel}
\setbeamerfont{frametitle}{size=\large,series=\bfseries}
\setbeamerfont{normal text}{size=\small}
\setbeamertemplate{frametitle}{\vskip0.12cm\insertframetitle\par\vskip0.08cm}
\setbeamertemplate{footline}{%
  \hbox{\begin{beamercolorbox}[wd=\paperwidth,ht=2.3ex,dp=0.9ex,
    leftskip=0.55cm,rightskip=0.55cm]{author in head/foot}%
    \scriptsize\color{muted}PyQCD / donghx\hfill
    4150 复现与对照\hfill\insertframenumber/\inserttotalframenumber
  \end{beamercolorbox}}}

\newcommand{\source}[1]{\vspace{0.05em}\par{\raggedright\scriptsize\color{muted}来源：#1\par}}
\newcommand{\codepath}[1]{\nolinkurl{#1}}
\newcommand{\verified}{\textcolor{evidence}{确证}}
\newcommand{\inferred}{\textcolor{accent}{推断}}
\newcommand{\unverified}{\textcolor{warning}{未验证}}
\newcolumntype{Y}{>{\raggedright\arraybackslash}X}
\newcolumntype{L}{>{\raggedright\arraybackslash}p{0.20\textwidth}}
\newcolumntype{M}{>{\raggedright\arraybackslash}p{0.28\textwidth}}

\tikzset{
  flow/.style={draw=accent,fill=accent!8,rounded corners=2pt,align=center,
    minimum height=0.78cm,text width=2.15cm,font=\small},
  state/.style={draw=primary,fill=primary!7,rounded corners=2pt,align=center,
    minimum height=0.72cm,text width=2.55cm,font=\small},
  arrow/.style={-{Latex[length=2mm]},draw=primary,semithick},
  labelstyle/.style={font=\scriptsize,inner sep=1.5pt,fill=white,text=muted}
}

\lstset{
  basicstyle=\ttfamily\scriptsize,backgroundcolor=\color{panel},frame=single,
  breaklines=true,breakatwhitespace=false,keepspaces=true,columns=flexible,
  numbers=left,numberstyle=\tiny\color{muted},keywordstyle=\color{primary}\bfseries,
  commentstyle=\color{evidence!70!black},stringstyle=\color{accent}
}

\title[4150 donghx 复现]{donghx 格点 QCD 算法的 PyQCD 复现}
\subtitle{组态 4150：eigvec / VdV / VVV / 2pt / OPE 的真实对照与下游边界}
\author[Codex]{Codex / PyQCD}
\institute{格点 QCD：蒸馏、质子相关函数与胶子算符}
\date{2026 年 8 月 29 日}

\begin{document}

\begin{frame}[plain]
  \titlepage
  \vfill
  \begin{center}
    \small 工作目录：\codepath{/root/PyQCD}\quad
    参考：\codepath{/root/PyQCD/refer/donghx}\quad
    单测组态：\texttt{4150}
  \end{center}
\end{frame}

\begin{frame}[t]{结论先行：低层七项与两种 2pt 投影均通过，完整链仍有明确边界}
  \begin{columns}[T,totalwidth=\textwidth]
    \begin{column}{0.30\textwidth}
      \begin{block}{低层对象}
        \centering\Large\textcolor{evidence}{7/7}\par\normalsize
        eigvec、phase、VdV、VVV、Clover、dual、OPE
      \end{block}
      \small 相对差为 $0$ 或 $1.4\times10^{-15}$ 量级。
    \end{column}
    \begin{column}{0.30\textwidth}
      \begin{block}{费米子 2pt}
        \centering\Large\textcolor{evidence}{2\ \text{投影}}\par\normalsize
        $Cg5g4$ 与 $Cg5$，$P=(0,0,0)$
      \end{block}
      \small 各计算 $t_{\rm source}=0$ 的 $\Delta t=2\ldots36$。
    \end{column}
    \begin{column}{0.30\textwidth}
      \begin{block}{证据边界}
        \centering\Large\textcolor{warning}{未验证}\par\normalsize
        外部逐时间 VVV、3pt、ratio、barematrix
      \end{block}
      \small 参考路径缺失或几何/输入不能一一配对，未以空数组代替结果。
    \end{column}
  \end{columns}
  \vfill
  \colorbox{good}{\parbox{0.94\linewidth}{\centering
    \verified\ 生产侧的 donghx 两项质子缩并已补入 \codepath{pyqcd/contraction/_donghx.py}；
    \inferred\ 复现结论限于当前可读输入和明确时间对。}}
  \source{实测：\codepath{examples/pyqcd/cmp1/v202608291500_lowlevel/results.json}；
  \codepath{examples/pyqcd/cmp1/v202608291710_fermion_cg5g4/results.json}；
  \codepath{logs/donghx_4150_20260829/summary.md}}
\end{frame}

\begin{frame}[t]{任务与输入：同一组态驱动参考公式和 PyQCD 生产实现}
  \footnotesize
  \setlength{\tabcolsep}{3pt}
  \renewcommand{\arraystretch}{0.98}
  \begin{tabularx}{0.96\textwidth}{@{}L Y M Y@{}}
    \toprule
    对象 & 4150 输入事实 & 参考结果/状态 & 验收方式\\
    \midrule
    Clover 规范场 & raw，573439152 B & Operator.py 可读 & 两时间片全张量\\
    eigvec & 72 个时间片，约 4.78 GB & Calc\_VVV.py 同式读取 & f8 交错复数、形状\\
    light peram & 288 个文件，约 13.27 GB & 用户给定路径存在 & $4\times4\times100\times100$\\
    HYP 3D(1/3/5)、4D(10) & 各 7 个 LIME record & 存在；本轮未重算 & 资产盘点\\
    2pt Result & 322 个文件，含 Cg5/Cg5g4 & 4150 可配对 & 选定时间对相对差\\
    TMD OPE Result & x/y/z 共 58 个文件；有横向位移 & 与直线 OPE 分开 & 不混合比较\\
    \bottomrule
  \end{tabularx}
  \vspace{0.12em}
  \begin{block}{判据}
    形状、轴序、dtype 和不变量先行；参考成品不存在即记为 \unverified，不把“能运行”升级为“一致”。
  \end{block}
  \source{盘点：\codepath{cmp1/v202608291424/manifest.json}；输入：\codepath{cmp1/manifest_4150.py}}
\end{frame}

\begin{frame}[t]{算法总览：蒸馏费米子链与胶子 OPE 链在 2pt/3pt 处汇合}
  \centering
  \begin{tikzpicture}[node distance=0.34cm and 0.20cm]
    \node[flow] (gauge) {规范场\ $U_\mu(x)$};
    \node[flow,right=of gauge] (clover) {Clover\ $F,\widetilde F$};
    \node[flow,right=of clover] (ope) {Wilson 线\ OPE $O(z)$};
    \node[flow,below=0.62cm of gauge] (eig) {低模\ $\phi_n(x)$};
    \node[flow,right=of eig] (vertex) {相位\ VdV / VVV};
    \node[flow,right=of vertex] (peram) {peram\ $\tau$};
    \node[flow,right=of peram] (c2) {Wick\ $C_2$};
    \node[flow,right=of c2] (analysis) {3pt / ratio\ / PDF};
    \draw[arrow] (gauge) -- (clover);
    \draw[arrow] (clover) -- (ope);
    \draw[arrow] (eig) -- (vertex);
    \draw[arrow] (vertex) -- (peram);
    \draw[arrow] (peram) -- (c2);
    \draw[arrow] (c2) -- (analysis);
    \draw[arrow] (ope) |- (analysis);
  \end{tikzpicture}
  \vfill
  \begin{tabularx}{0.92\textwidth}{@{}lY@{}}
    \toprule
    物理分支 & 输入、状态、输出\\
    \midrule
    蒸馏质子 & $\phi_n$ 低模基 $\to$ 动量投影顶点 $\to$ perambulator $\to$ 颜色 epsilon 两项缩并 $\to C_2$。\\
    胶子算符 & 原始 link $\to$ Clover 场强/对偶场强 $\to$ 直线 Wilson 线 $\to O(z)$；横向 staple 属于另一个几何。\\
    下游统计 & $C_3/C_2$、真空扣除、拟合和 bare matrix 需要可配对的 3pt/样本输入，本轮不越过缺口。\\
    \bottomrule
  \end{tabularx}
  \source{参考流程：\codepath{refer/donghx/Calc_VVV.py:45--136}、\codepath{refer/donghx/Operator.py:164--184,339--368}；PyQCD 接口：\codepath{pyqcd/vertex/_vertex.py}、\codepath{pyqcd/operator/_gluon_ope.py}}
\end{frame}

\begin{frame}[t]{eigvec、相位与 VdV/VVV：颜色结构决定可比较的轴序}
  \footnotesize
  \begin{align*}
    \phi_n(x)&=r_{n}(x)+i\,s_{n}(x),
    &e_p(x)&=\exp\!\left[-\frac{2\pi i}{N_x}(p_z z+p_y y+p_x x)\right],\\
    V_{mn}(p)&=\sum_{x,a}\phi_m^{a*}(x)e_p(x)\phi_n^a(x),
    &V_{mnl}(p)&=\sum_x e_p(x)\,\epsilon_{abc}\phi_m^a(x)\phi_n^b(x)\phi_l^c(x).
  \end{align*}
  \begin{tabularx}{0.96\textwidth}{@{}L Y@{}}
    \toprule
    状态/动作 & 实现与检查\\
    \midrule
    二进制读取 & f8 交错的实虚部先 reshape 为 $(N_{\rm ev},N_x,N_x,N_x,3,2)$，再组成 complex；不把时间轴混入低模轴。\\
    动量投影 & VdV 的 phase 复制到颜色轴；VVV 使用标量 phase 和三色 epsilon，保持 $(m,n,l)$ 的六置换符号。\\
    极端检查 & $p=0$ 时 phase 恒为 1；重复低模时 VVV 的完全反对称性使重复指标分量消失。\\
    \bottomrule
  \end{tabularx}
  \vfill
  \begin{block}{本轮证据}
    真实 4150：eigvec $(8,13824,3)$、phase $(13824,)$、VdV $(1,4,4)$、VVV $(4,4,4)$；
    相对差为 $0,0,1.31\times10^{-15},4.75\times10^{-16}$。
  \end{block}
  \source{参考：\codepath{Calc_VVV.py:28--136}；实现：\codepath{tools/_io.py}、\codepath{vertex/_vertex.py}；结果：\codepath{lowlevel/results.json}}
\end{frame}

\begin{frame}[t]{Clover 与对偶场强（1/2）：先固定局域张量的归一化}
  \footnotesize
  \begin{align*}
    F_{\mu\nu}(x)&=-\frac{i}{8}\bigl(C_{\mu\nu}(x)-C_{\mu\nu}^{\dagger}(x)\bigr),
    &\widetilde F_{\mu\nu}(x)&=\frac12\epsilon_{\mu\nu\rho\sigma}F_{\rho\sigma}(x).
  \end{align*}
  \begin{tabularx}{0.96\textwidth}{@{}L Y@{}}
    \toprule
    步骤 & 物理/工程约束\\
    \midrule
    Clover & 四个方向的正反 plaquette 组成反厄米场强；颜色矩阵最后两轴保持 $(3,3)$。\\
    对偶 & epsilon 收缩只使用明确的 $(\mu,\nu,\rho,\sigma)$ 约定；本轮比较包含参考固定轴序候选。\\
    \bottomrule
  \end{tabularx}
  \vspace{0.12em}
  \begin{block}{真实比较}
    两个时间片的 Clover / dual 相对差分别为 $2.45\times10^{-16}$、$1.89\times10^{-16}$。
  \end{block}
  \source{参考：\codepath{Operator.py:150--184}；实现：\codepath{operator/_gluon_ope.py:63--133}；结果：\codepath{lowlevel/results.json}}
\end{frame}

\begin{frame}[t]{直线 Wilson 线 OPE（2/2）：几何与正负 $z$ 分支必须独立}
  \footnotesize
  \begin{align*}
    M_{\mu\nu;\rho\sigma}(z)&=\sum_{\boldsymbol x}\operatorname{Tr}\!\left[
      F_{\mu\nu}(x+z\hat z)W_{0\to z}^{\dagger}\widetilde F_{\rho\sigma}(x)W_{0\to z}\right],\\
    O(z)&=M_{tx;tx}(z)+M_{ty;ty}(z)-2M_{xy;xy}(z).
  \end{align*}
  \begin{tabularx}{0.96\textwidth}{@{}L Y@{}}
    \toprule
    步骤 & 物理/工程约束\\
    \midrule
    $+z$ Wilson 线 & 逐步左乘 $U^\dagger$，插入第二场强，再右乘 $U$；输出按时间片求空间和。\\
    $-z$ 变体 & 使用独立的反向 link 链；方向不能用取绝对值或取实部来修正。\\
    几何边界 & 本轮是直线 $z$，返回 $(\Delta z,N_t)$；参考 TMD 的横向 $\Delta x/\Delta y$ staple 另属一类算符。\\
    \bottomrule
  \end{tabularx}
  \vspace{0.12em}
  \begin{block}{真实比较}
    $z=0,1$、两个时间片的直线 OPE 相对差为 $1.36\times10^{-15}$；该数值不代表横向 staple 已验证。
  \end{block}
  \source{参考：\codepath{Operator.py:339--399}；实现：\codepath{operator/_gluon_ope.py:135--221}；结果：\codepath{lowlevel/results.json}}
\end{frame}

\begin{frame}[t]{质子 2pt：peram 轴序与两项颜色 epsilon 缩并被显式复现}
  \footnotesize
  对固定 $(t_s,t_0)$，令 $\tau_{d_s d_0;\,e_s e_0}$ 为 $(4,4,N_{\rm ev},N_{\rm ev})$ peram，
  $I_1,I_2$ 为插值矩阵，并定义
  \[
    A_{g j b e}=\sum_{h k}(I_1)_{g h}\,\tau_{h k;b e}\,(I_2)_{j k}.
  \]
  donghx 的质子矩阵为
  \[
    C_{il}=\epsilon_{abc}\,\tau_{g j;a d}\,A_{g j;b e}\,
      \tau_{i l;c f}\,V^{*}_{def}
    -\epsilon_{abc}\,\tau_{g l;a f}\,A_{g j;b e}\,
      \tau_{i j;c d}\,V^{*}_{def}.
  \]
  \begin{tabularx}{0.96\textwidth}{@{}L Y@{}}
    \toprule
    参考动作 & PyQCD 复现\\
    \midrule
    读取 & 四个 source-Dirac 文件合并，再 reshape/transponse 为 $(t_{\rm sink},d_{\rm sink},d_{\rm source},e_s,e_0)$。\\
    插值 & 支持 $Cg5$、$Cg5g3$、$Cg5g4$ 与三个 offdiag 变体；主对照为 $Cg5g4$。\\
    缩并 & 生产接口 \codepath{contract_donghx_2pt_pair} 保留上述两项，不复用位置不同的通用 Wick gamma 约定。\\
    轴检验 & peram 必须为 $(4,4,N_{\rm ev},N_{\rm ev})$，两个 VVV 必须为 $(N_{\rm ev})^3$；不匹配立即报错。\\
    \bottomrule
  \end{tabularx}
  \source{参考缩并：\codepath{refer/donghx/2pt_proton_cupy_mom0_Liu_Cg5gmu_L24x72_new_cpu.py:346--381}；PyQCD：\codepath{pyqcd/contraction/_donghx.py:85--151}}
\end{frame}

\begin{frame}[t]{投影与反周期边界：时间半平面的符号不能在统计后补救}
  \footnotesize
  \begin{align*}
    P_+&=\frac12(\gamma_0+\gamma_4), & P_-&=\frac12(\gamma_0-\gamma_4),\\
    C_+(t_s,t_0)&=\operatorname{tr}[P_+C(t_s,t_0)],
    &C_-(t_s,t_0)&=\operatorname{tr}[P_-C(t_s,t_0)].
  \end{align*}
  \renewcommand{\arraystretch}{1.00}
  \begin{tabularx}{0.96\textwidth}{@{}L Y@{}}
    \toprule
    时间关系 & 反周期边界处理\\
    \midrule
    $t_s<t_0$ & $P_+$ 通道乘 $-1$（反向传播跨越时间边界）。\\
    $t_s>t_0$ & $P_-$ 通道乘 $-1$；本轮 $t_0=0$，选定 $t_s=2\ldots36$。\\
    未计算的时间对 & 保持零；只抽取相同 $(t_s,t_0)$，不比较稀疏数组与完整矩阵。\\
    物理检查 & $p=0$、投影和边界符号均受控测试；不以 $|C|$ 或 $\operatorname{Re}C$ 掩盖错误。\\
    \bottomrule
  \end{tabularx}
  \colorbox{warnbg}{\parbox{0.94\linewidth}{\centering
    参考保存完整 $(72,72,4,4)$；本轮只算 $t_0=0$ 的 35 对，故不等于全 source 时间完成。}}
  \source{参考：\codepath{2pt_proton...new_cpu.py:418--465}；实现：\codepath{contraction/_baroperator.py:332--365}}
\end{frame}

\begin{frame}[t]{低层真实结果：七个对象在同一 4150 规范场上闭合}
  \footnotesize
  \setlength{\tabcolsep}{4pt}
  \renewcommand{\arraystretch}{1.10}
  \begin{tabularx}{0.96\textwidth}{@{}L Y r r r l@{}}
    \toprule
    对象 & 结果形状 & 参考秒 & PyQCD 秒 & 相对差 & 状态\\
    \midrule
    eigvec & $(8,13824,3)$ & 0.0765 & 0.1068 & $0$ & \verified\\
    phase & $(13824)$ & 0.0980 & 1.0491 & $0$ & \verified\\
    VdV & $(1,4,4)$ & 0.1536 & 0.0029 & $1.31\times10^{-15}$ & \verified\\
    VVV & $(4,4,4)$ & 0.1075 & 0.0141 & $4.75\times10^{-16}$ & \verified\\
    Clover & $(4,4,2,24,24,24,3,3)$ & 1.9927 & 1.0342 & $2.45\times10^{-16}$ & \verified\\
    dual & $(4,4,2,24,24,24,3,3)$ & 0.0681 & 1.0365 & $1.89\times10^{-16}$ & \verified\\
    OPE & $(2,2)$ & 0.0973 & 0.2419 & $1.36\times10^{-15}$ & \verified\\
    \bottomrule
  \end{tabularx}
  \vfill
  \begin{block}{解读}
    低层差异均低于各自 $10^{-9}$--$10^{-14}$ 容差；时间是 CPU 单次函数计时，不能直接当作跨后端性能结论。
    OPE 这里对照的是同一规范场上的参考公式重写，外部 TMD staple 成品另行标记。
  \end{block}
  \source{\codepath{examples/pyqcd/cmp1/v202608291500_lowlevel/results.json}；命令：\texttt{python}\,\codepath{examples/pyqcd/cmp1/run_4150_lowlevel.py --conf 4150 --outdir ...}}
\end{frame}

\begin{frame}[t]{2pt 真实结果：Cg5g4 与 Cg5 都与可配对的 4150 成品一致}
  \footnotesize
  \setlength{\tabcolsep}{2.5pt}
  \renewcommand{\arraystretch}{1.02}
  \begin{tabularx}{0.96\textwidth}{@{}Y Y r r Y@{}}
    \toprule
    通道/动量 & 计算范围 & contract 相对差 & nopol 相对差 & 最大 $|\Delta|$；参考 dtype\\
    \midrule
    $Cg5g4$, $P=0$ & 35 对；$t_0=0$ & $3.50\times10^{-15}$ & $1.27\times10^{-15}$ & $4.25\times10^{-17}$；c128\\
    $Cg5$, $P=0$ & 35 对；$t_0=0$ & $2.29\times10^{-15}$ & $1.53\times10^{-15}$ & $8.33\times10^{-17}$；c128\\
    $Cg5g4$, $P_z=1$ & $\Delta t=2$ & $2.53\times10^{-15}$ & $8.56\times10^{-16}$ & $1.73\times10^{-17}$；c128\\
    $Cg5g4$, $P_z=2$ & $\Delta t=2$ & $2.31\times10^{-15}$ & $8.26\times10^{-16}$ & $1.30\times10^{-17}$；c128\\
    $Cg5g4$, $P_z=5$ & $\Delta t=2$ & $1.66\times10^{-15}$ & $1.42\times10^{-16}$ & $2.63\times10^{-20}$；c128\\
    $Cg5g4$, $P_x=1$ & $\Delta t=2$ & $2.64\times10^{-6}$ & $1.42\times10^{-6}$ & $2.32\times10^{-8}$；c64，$\mathrm{tol}=10^{-5}$\\
    $Cg5g4$, $P_y=1$ & $\Delta t=2$ & $2.59\times10^{-6}$ & $5.94\times10^{-7}$ & $1.90\times10^{-8}$；c64，$\mathrm{tol}=10^{-5}$\\
    \bottomrule
  \end{tabularx}
  \vspace{0.10em}
  \begin{block}{关键状态}
    contract/nopol 均为 \verified；逐时间 VVV 参考缺失，故 runner 整体按最严重项记为 \unverified。
  \end{block}
  \source{主结果：\codepath{cmp1/v202608291710_fermion_cg5g4/results.json}；Cg5：\codepath{cmp1/v202608291705_fermion_cg5/results.json}；索引：\codepath{logs/donghx_4150_20260829/summary.md}}
\end{frame}

\begin{frame}[t]{smear 与动量边界：本轮能闭合的范围和不能越过的输入缺口}
  \footnotesize
  \renewcommand{\arraystretch}{1.00}
  \begin{tabularx}{0.96\textwidth}{@{}L Y Y@{}}
    \toprule
    分支 & 当前事实 & 判定\\
    \midrule
    无 momentum smear & 标准 light peram + eigvec；$Cg5/Cg5g4$、$P=0$ 的 35 对通过。 & \verified\\
    无 momentum smear 的方向/大小 & $P_z=1,2,5$、$P_x=1$、$P_y=1$；按参考 dtype 给容差。 & \verified\\
    momentum smear $\pm2x/\pm2y/\pm2z$ & 参考有结果，但需独立 smeared peram；本轮只有标准 light peram。 & \unverified\\
    gauge smear & HYP record 存在；尚未完成 record $\to$ PyQCD link $\to$ Clover/OPE。 & \unverified\\
    \bottomrule
  \end{tabularx}
  \begin{block}{物理含义}
    smearing 改变输入链：momentum 同时改变 phase/peram；gauge 必须在同一 HYP link 上重算 Clover/OPE。
  \end{block}
  \source{参考：\codepath{2pt_proton...mom2_xdir_dcu.py:57--180}；资产：\codepath{cmp1/v202608291424/manifest.json}}
\end{frame}

\begin{frame}[t]{3pt、ratio、barematrix：接口可审阅，但缺少可配对的 4150 样本}
  \footnotesize
  \renewcommand{\arraystretch}{1.00}
  \begin{tabularx}{0.96\textwidth}{@{}L Y Y@{}}
    \toprule
    层级 & PyQCD 入口/物理定义 & 本轮状态\\
    \midrule
    3pt & 顺序源/插入时间的三点缩并，再与 2pt 做真空扣除或 ratio。 & \unverified：未找到可配对成品。\\
    ratio & \codepath{analysis/_analyse.py:361} 的 $C_3/C_2$，以及 \codepath{analysis/_ratio2pt.py:387} 的逐 $z$ 拟合。 & 接口可审阅；无 4150 原始 3pt 样本。\\
    barematrix & \codepath{analysis/_bare_matrix.py:132} 按方向组织 ratio 和拟合。 & \unverified：缺少方向样本。\\
    \bottomrule
  \end{tabularx}
  \vspace{0.12em}
  \colorbox{warnbg}{\parbox{0.94\linewidth}{\centering
    最短闭环是“可配对 3pt/peram + 参考 ratio/barematrix 输出”；否则几何/样本混合，新增数值不可比。}}
  \source{实现：\codepath{analysis/_analyse.py}、\codepath{analysis/_ratio2pt.py}、\codepath{analysis/_bare_matrix.py}；缺口：\codepath{cmp1/v202608291424/manifest.json}}
\end{frame}

\begin{frame}[t]{TMD 几何边界：直线 OPE 通过不等于横向 staple 闭环}
  \footnotesize
  \begin{tabularx}{0.96\textwidth}{@{}L Y Y@{}}
    \toprule
    层级 & 本轮可追溯事实 & 判定\\
    \midrule
    直线低层 & $O(z)$ 的 Clover、对偶场强、Wilson 线和两个 $z$ 点已对照；相对差 $1.36\times10^{-15}$。 & \verified\\
    参考 TMD 几何 & 成品含横向 $\Delta x/\Delta y$ staple；不是本轮直线 $z$ 算符的同一输入。 & \unverified\\
    完整 TMD-PDF & 还需 $b_\perp$、soft/rapidity 减除、匹配及 $a/P_z/m_\pi/L$ 多尺度外推。 & \unverified\\
    \bottomrule
  \end{tabularx}
  \vspace{0.12em}
  \begin{block}{物理判断}
    只把同一 gauge、同一几何、同一样本的结果放入 ratio；本轮证据停在直线 OPE 低层，不宣称完整 TMD-PDF。
  \end{block}
  \source{直线结果：\codepath{cmp1/v202608291500_lowlevel/results.json}；参考 TMD 资产：\codepath{cmp1/v202608291424/manifest.json}}
\end{frame}

\begin{frame}[t]{本轮 debug 与优化：差异来自输入/命名/证据等级，而非数值算法漂移}
  \footnotesize
  \renewcommand{\arraystretch}{1.00}
  \begin{tabularx}{0.96\textwidth}{@{}L Y Y@{}}
    \toprule
    发现 & 处理 & 证据\\
    \midrule
    VVV 缓存 & 校验 shape/dtype 后复用 $(N_t,N_{\rm ev}^3)$ 的 c128 缓存。 & runner 7/7；日志有 reuse。\\
    Cg5 命名 & 为不带 \codepath{_Cg5} 后缀的参考文件增加显式分支。 & Cg5 35 对通过。\\
    参考 VVV 缺失 & 不伪造比较；记录为 ref 缺失并整体降级。 & 状态正确为 \unverified。\\
    复数精度 & c128 逐位；Px/Py 为 c64，采用 $10^{-5}$ 容差。 & 单时间对均在容差内。\\
    \bottomrule
  \end{tabularx}
  \begin{block}{保持的物理不变量}
    不随意取实部或绝对值；VVV 保持六置换；peram/VVV 轴不隐式转置。
  \end{block}
  \source{实现：\codepath{cmp1/run_4150_fermion.py:98--180}、\codepath{cmp1/cases_4150_fermion.py:46--71}；测试：\codepath{cmp1/verify_4150_fermion_runner.py}}
\end{frame}

\begin{frame}[t]{验收与复现：命令、结果文件和状态语义均已留痕}
  \footnotesize
  \renewcommand{\arraystretch}{1.00}
  \begin{tabularx}{0.96\textwidth}{@{}L Y l@{}}
    \toprule
    层级 & 命令/产物 & 结果\\
    \midrule
    资产 & \codepath{verify_manifest_4150.py} & PASS 4/4\\
    低层受控契约 & \codepath{verify_4150_lowlevel.py} & PASS 3/3\\
    2pt 受控契约 & \codepath{verify_4150_fermion.py} & PASS 3/3\\
    runner 契约 & \codepath{verify_4150_fermion_runner.py} & PASS 7/7\\
    低层真实 & \codepath{v202608291500_lowlevel/results.json} & 7 pass\\
    2pt 真实 & \codepath{v202608291710_fermion_cg5g4/results.json} & 2 pass + 1 unverified\\
    \bottomrule
  \end{tabularx}
  \vspace{0.10em}
  \begin{block}{主回归}
    \texttt{python}\,\codepath{examples/pyqcd/conftest.py}：41 passed，0 failed；编译验收仍要求两遍 XeLaTeX、零溢出且页数闭合。
  \end{block}
  \source{证据索引：\codepath{logs/donghx_4150_20260829/summary.md}；产物：\codepath{examples/pyqcd/cmp1/}}
\end{frame}

\appendix

\begin{frame}[t]{附录：参考代码到 PyQCD 的逐层映射（1/2）}
  \footnotesize
  \setlength{\tabcolsep}{3pt}
  \renewcommand{\arraystretch}{1.00}
  \begin{tabularx}{0.96\textwidth}{@{}L Y Y@{}}
    \toprule
    参考对象 & donghx 关键位置 & PyQCD 对应位置\\
    \midrule
    eigvec/phase & \codepath{Calc_VVV.py:28--59} & \codepath{tools/_io.py:32--64}；\codepath{vertex/_vertex.py:28--106}\\
    VdV/VVV & \codepath{Calc_VVV.py:62--136} & \codepath{vertex/_vertex.py:113--207}\\
    Clover/dual & \codepath{Operator.py:150--184} & \codepath{operator/_gluon_ope.py:63--133}\\
    peram 轴序 & \codepath{2pt_proton...new_cpu.py:193--224} & \codepath{tools/_io.py:115--223}\\
    2pt 缩并 & \codepath{2pt_proton...new_cpu.py:346--381} & \codepath{contraction/_donghx.py:85--151}\\
    投影/边界 & \codepath{2pt_proton...new_cpu.py:418--465} & \codepath{contraction/_baroperator.py:332--365}\\
    \bottomrule
  \end{tabularx}
  \source{参考：\codepath{refer/donghx}；实现：\codepath{pyqcd/}；本表列出本轮直接使用的入口}
\end{frame}

\begin{frame}[t]{附录：参考代码到 PyQCD 的逐层映射（2/2）}
  \footnotesize
  \setlength{\tabcolsep}{3pt}
  \renewcommand{\arraystretch}{1.00}
  \begin{tabularx}{0.96\textwidth}{@{}L Y Y@{}}
    \toprule
    参考对象 & donghx 关键位置 & PyQCD 对应位置\\
    \midrule
    直线 OPE & \codepath{Operator.py:339--399} & \codepath{operator/_gluon_ope.py:135--221}\\
    下游统计 & 参考脚本/结果缺少可配对样本 & \codepath{analysis/_analyse.py:361,744}；\codepath{analysis/_bare_matrix.py:132}\\
    TMD 几何 & 结果含横向 staple，不能与直线 $z$ 混合 & \codepath{renorm/_tmd.py} 等待同几何输入\\
    本轮证据 & 低层 results.json；2pt results.json & 状态字段区分 \verified、\unverified\\
    \bottomrule
  \end{tabularx}
  \vspace{0.12em}
  \begin{block}{追溯规则}
    先固定参考代码入口、输入几何和 dtype，再解释相对差；缺失成品只保留接口/资产事实，不填补数值。
  \end{block}
  \source{结果索引：\codepath{logs/donghx_4150_20260829/summary.md}；完整代码在 \codepath{refer/donghx} 与 \codepath{pyqcd/}}
\end{frame}

\begin{frame}[t]{附录：下一步的可验收闭环（1/2）：先补齐 2pt 输入与 VVV 成品}
  \footnotesize
  \renewcommand{\arraystretch}{1.00}
  \begin{tabularx}{0.96\textwidth}{@{}L Y Y Y@{}}
    \toprule
    任务 & 前置输入 & 产物/判据 & 当前阻塞\\
    \midrule
    2pt 变体矩阵 & momentum-smeared peram；HYP link & 各 smear/方向/大小独立 JSON；contract/nopol 在 dtype 容差内 & 未给出 smeared peram\\
    VVV 成品闭环 & 参考 \codepath{VVV.t*.Px*Py*Pz*} 或等价文件 & 逐时间 VVV 相对差；状态不再 unverified & 结果目录未见该文件\\
    \bottomrule
  \end{tabularx}
  \vspace{0.12em}
  \begin{block}{交付判据}
    新输入到位后复用 manifest/runner；每层保存中间数组、最终数组、dtype、形状、耗时和容差。
  \end{block}
  \source{计划：\codepath{docs/superpowers/plans/2026-08-29-donghx-4150-reproduction.md}；边界：\codepath{logs/donghx_4150_20260829/summary.md}}
\end{frame}

\begin{frame}[t]{附录：下一步的可验收闭环（2/2）：再扩展 3pt、HYP 与完整 TMD}
  \footnotesize
  \renewcommand{\arraystretch}{1.00}
  \begin{tabularx}{0.96\textwidth}{@{}L Y Y Y@{}}
    \toprule
    任务 & 前置输入 & 产物/判据 & 当前阻塞\\
    \midrule
    3pt/ratio & 顺序源、插入时间、OPE 样本 & raw $C_3/C_2$、真空扣除、拟合表、barematrix & 外部目录无法配对\\
    HYP/OPE & record 可读的 link 适配 & 各 3D/4D smear 的 Clover、OPE、几何元数据 & 本轮只盘点资产\\
    TMD & $z,b_\perp,\tau$、soft/rapidity/matching、多尺度数据 & 软因子、匹配、连续极限闭环 & 进阶项，当前不宣称完成\\
    \bottomrule
  \end{tabularx}
  \vspace{0.12em}
  \begin{block}{物理边界}
    只有同一 link 几何、同一样本和同一归一化链闭合后，才进入最终 PDF/拟合结论。
  \end{block}
  \source{当前边界：\codepath{logs/donghx_4150_20260829/summary.md}；下游入口：\codepath{pyqcd/analysis/}、\codepath{pyqcd/renorm/}}
\end{frame}

\end{document}
